\section{Sequential competition trace}

\begin{table}[t]
  \caption{Sequential V33 live results. Ratings are displayed competition
  values; stopping is adaptive, so changes are descriptive.}
  \label{tab:v33-results}
  \centering
  \small
  \begin{tabular}{lrrrrl}
    \toprule
    Family & Games & Start & Final & Change & Stop \\
    \midrule
    Bargaining & 150 & 1344.98 & 1535.92 & $+190.94$ & target \\
    Persuasion & 78 & 1770.17 & 1825.84 & $+55.67$ & drawdown \\
    Negotiation & 20 & 1695.03 & 1661.30 & $-33.73$ & drawdown \\
    \bottomrule
  \end{tabular}
\end{table}


The August 30, 2026 run targeted 150 authoritative completions per family in
fixed order: bargaining, persuasion, then negotiation. Matchmaking used
concurrency one, and adjacent dashboard snapshots supplied assignment-level
displayed rating changes. Bargaining reached 150 games. Persuasion stopped at
78 after a 23.19-point drawdown from its session peak; negotiation stopped at
20 after losing 33.73 from its initial and peak rating.

Role sums were Alice $+80.15$ ($n=64$), Bob $+110.79$ ($n=86$), persuasion
buyer $+34.77$ ($n=37$), persuasion seller $+20.90$ ($n=41$), negotiation buyer
$-20.33$ ($n=11$), and negotiation seller $-13.40$ ($n=9$). Bargaining had
76/1/73 positive/zero/negative changes; persuasion 38/0/40; negotiation 7/0/13.
The positive totals therefore do not imply that most assignments gained rating.

We do not compare marginal persuasion information regimes. All 26 hidden-value
games assigned V33 the seller role, including all 15 hidden-value text games;
there is no hidden-value buyer support. The previously tempting comparison with
known-value text games therefore confounds role and information. Appendix~
\ref{app:statistics} reports the complete support table and descriptive paths.

The trace has no contemporaneous baseline or randomized branch allocation.
Opponent repetition is unavailable, ratings and population may drift, policies
were selected from prior live runs, and stopping depends on outcomes. We report
no inferential interval or mechanism-effect estimate. A valid follow-up should
interleave frozen branches within role--configuration strata, retain privacy-safe
opponent identifiers, preregister guards, and use cluster-aware anytime-valid
inference. Offline tests passed and no live action fallback or assignment
overshoot occurred; a telemetry-order bug is documented in the supplement.
